% Generated by supplementary/build_additional_file_1.py from the accepted tables. Each table is
% stored by \supptable{number}{...} and printed in numeric order by additional_file_1.tex.
\supptable{2}{%
\begin{table}[H]
\centering\footnotesize
\caption{\textbf{Supplemental Table S2. Concordance with Space Ranger for every resolution policy, scale and slide.} Mass ratio is the STARgrid molecule total divided by the Space Ranger total. Exact entries are the fraction of entries nonzero in either matrix that have identical counts. Correlations are Pearson $r$. SPATCH uses the bins in Space Ranger's published filtered matrix at 8\,\textmu m.}
\label{supptab:2}
\begin{tabular}{llllrrr}
\toprule
Slide & Bin (\textmu m) & Policy & Mass ratio & Exact entries & Gene-total $r$ & Bin-total $r$ \\
\midrule
Colorectal & 2 & \policyStrict{} & 0.8815 & 0.8109 & 0.999998 & 0.98046 \\
Colorectal & 2 & \policySoft{} & 1.0454 & 0.6589 & 0.999999 & 0.98373 \\
Colorectal & 2 & \policyHard{} & 1.0462 & 0.8456 & 0.999999 & 0.98430 \\
Colorectal & 2 & \policyGated{} & 0.9515 & 0.8197 & 0.999999 & 0.97900 \\
Colorectal & 8 & \policyStrict{} & 0.8815 & 0.8091 & 0.999998 & 0.99848 \\
Colorectal & 8 & \policySoft{} & 1.0454 & 0.8015 & 0.999999 & 0.99910 \\
Colorectal & 8 & \policyHard{} & 1.0462 & 0.9045 & 0.999999 & 0.99926 \\
Colorectal & 8 & \policyGated{} & 0.9515 & 0.8548 & 0.999999 & 0.99877 \\
Colorectal & 16 & \policyStrict{} & 0.8815 & 0.7993 & 0.999998 & 0.99950 \\
Colorectal & 16 & \policySoft{} & 1.0454 & 0.8484 & 0.999999 & 0.99975 \\
Colorectal & 16 & \policyHard{} & 1.0462 & 0.9148 & 0.999999 & 0.99980 \\
Colorectal & 16 & \policyGated{} & 0.9515 & 0.8551 & 0.999999 & 0.99964 \\
H1-VM2JXXK/A1 & 2 & \policyStrict{} & 0.7953 & 0.7462 & 0.9999981 & 0.97975 \\
H1-VM2JXXK/A1 & 2 & \policySoft{} & 1.0262 & 0.5581 & 0.9999997 & 0.98744 \\
H1-VM2JXXK/A1 & 2 & \policyHard{} & 1.0263 & 0.8850 & 0.9999997 & 0.98916 \\
H1-VM2JXXK/A1 & 2 & \policyGated{} & 0.8783 & 0.8063 & 0.9999959 & 0.98431 \\
H1-VM2JXXK/A1 & 8 & \policyStrict{} & 0.7953 & 0.7359 & 0.9999981 & 0.99802 \\
H1-VM2JXXK/A1 & 8 & \policySoft{} & 1.0262 & 0.7322 & 0.9999997 & 0.99901 \\
H1-VM2JXXK/A1 & 8 & \policyHard{} & 1.0263 & 0.9203 & 0.9999997 & 0.99935 \\
H1-VM2JXXK/A1 & 8 & \policyGated{} & 0.8783 & 0.8103 & 0.9999959 & 0.99853 \\
H1-VM2JXXK/A1 & 16 & \policyStrict{} & 0.7953 & 0.7072 & 0.9999981 & 0.99924 \\
H1-VM2JXXK/A1 & 16 & \policySoft{} & 1.0262 & 0.7914 & 0.9999997 & 0.99971 \\
H1-VM2JXXK/A1 & 16 & \policyHard{} & 1.0263 & 0.9237 & 0.9999997 & 0.99981 \\
H1-VM2JXXK/A1 & 16 & \policyGated{} & 0.8783 & 0.7905 & 0.9999959 & 0.99948 \\
Ovarian & 2 & \policyStrict{} & 0.8760 & 0.7633 & 0.998937 & 0.98186 \\
Ovarian & 2 & \policySoft{} & 1.0735 & 0.5881 & 0.998913 & 0.98429 \\
Ovarian & 2 & \policyHard{} & 1.0764 & 0.8028 & 0.998925 & 0.98608 \\
Ovarian & 2 & \policyGated{} & 0.9577 & 0.7748 & 0.998949 & 0.97943 \\
Ovarian & 8 & \policyStrict{} & 0.8760 & 0.7427 & 0.998937 & 0.99854 \\
Ovarian & 8 & \policySoft{} & 1.0735 & 0.7212 & 0.998913 & 0.99910 \\
Ovarian & 8 & \policyHard{} & 1.0764 & 0.8543 & 0.998925 & 0.99936 \\
Ovarian & 8 & \policyGated{} & 0.9577 & 0.7975 & 0.998949 & 0.99882 \\
Ovarian & 16 & \policyStrict{} & 0.8760 & 0.7214 & 0.998937 & 0.99943 \\
Ovarian & 16 & \policySoft{} & 1.0735 & 0.7660 & 0.998913 & 0.99971 \\
Ovarian & 16 & \policyHard{} & 1.0764 & 0.8564 & 0.998925 & 0.99979 \\
Ovarian & 16 & \policyGated{} & 0.9577 & 0.7878 & 0.998949 & 0.99961 \\
SPATCH & 8 & \policyStrict{} & 0.8054 & --- & 0.999995 & 0.99663 \\
SPATCH & 8 & \policySoft{} & 1.0229 & --- & 0.999997 & 0.99841 \\
SPATCH & 8 & \policyHard{} & 1.0232 & --- & 0.999997 & 0.99877 \\
SPATCH & 8 & \policyGated{} & 0.8876 & --- & 0.999996 & 0.99756 \\
\bottomrule
\end{tabular}
\end{table}
}

\supptable{14}{%
\begin{table}[H]
\centering\footnotesize
\caption{\textbf{Supplemental Table S14. Registration sensitivity on the colorectal slide.} Each entry gives the score with Space Ranger's registration of the microscope image to the CytAssist image, followed by the score with ours, using the same image masks placed through each (Supplemental Fig.~S4). The largest change in any STARgrid minus Space Ranger score difference is $1.97\times10^{-3}$.}
\label{supptab:14}
\begin{tabular}{lrrr}
\toprule
Counts & In-cell fraction & Cell ROC AUC & Nucleus ROC AUC \\
\midrule
Space Ranger & 0.9494 / 0.9483 & 0.8522 / 0.8511 & 0.6980 / 0.7048 \\
STARgrid \policyHard{} & 0.9491 / 0.9481 & 0.8549 / 0.8537 & 0.6996 / 0.7064 \\
\policyHard{}, STARgrid only & 0.9429 / 0.9417 & 0.6751 / 0.6743 & 0.6133 / 0.6172 \\
\policyHard{}, Space Ranger only & 0.9428 / 0.9416 & 0.5983 / 0.5978 & 0.5641 / 0.5655 \\
\bottomrule
\end{tabular}
\end{table}
}

\supptable{13}{%
\begin{table}[H]
\centering\footnotesize
\caption{\textbf{Supplemental Table S13. Sensitivity of the top-gene comparison with Xenium.} Ovarian dataset against the adjacent-section Xenium run. The first row is the primary comparison reported in the main text: the 20 cancer-panel genes with the most Xenium counts, on 128\,\textmu m patches, with Pearson correlation. Each other row changes one of these choices; the variants are reported without adjustment for multiple testing. Genes: the number of genes compared; only 47 panel genes are measured in all count matrices, so the 50-gene variant is the all-gene variant. \policyHard{} wins and losses: genes whose correlation with Xenium across patches is higher or lower under \policyHard{} than under Space Ranger; no gene was tied and no correlation was undefined. $P$: exact two-sided sign test on the wins and losses. Median correlation: the median over the genes of the correlation of each set of counts with Xenium across patches.}
\label{supptab:13}
\begin{tabular}{lrrlrrrrr}
\toprule
 & & & & & & & \multicolumn{2}{c}{Median correlation} \\
\cmidrule(l){8-9}
Variant & Genes & Patch (\textmu m) & Correlation & \policyHard{} wins & \policyHard{} losses & $P$ & Space Ranger & \policyHard{} \\
\midrule
\textbf{Primary} & 20 & 128 & Pearson & 15 & 5 & 0.0414 & 0.238 & 0.257 \\
Fewer genes & 10 & 128 & Pearson & 8 & 2 & 0.109 & 0.188 & 0.201 \\
More genes & 30 & 128 & Pearson & 21 & 9 & 0.0428 & 0.218 & 0.210 \\
All eligible genes & 47 & 128 & Pearson & 33 & 14 & 0.00794 & 0.143 & 0.154 \\
Smaller patches & 20 & 64 & Pearson & 13 & 7 & 0.263 & 0.109 & 0.111 \\
Larger patches & 20 & 256 & Pearson & 13 & 7 & 0.263 & 0.430 & 0.454 \\
Rank correlation & 20 & 128 & Spearman & 18 & 2 & 0.000402 & 0.238 & 0.247 \\
\bottomrule
\end{tabular}
\end{table}
}

\supptable{7}{%
\begin{table}[H]
\centering\scriptsize
\caption{\textbf{Supplemental Table S7. Nucleus by nucleus: paired depth in place and against a far-shift background.} For each nucleus with at least 4 capture locations in the tissue, the paired depth is one minus the mean count on its locations divided by the mean on the same locations moved to each of the 8 neighbouring shifts, one location away (neighbours), or to each of the 16 next shifts out, two locations away (next out). In place and Background are medians in \%, and In place $-$ background is their difference in percentage points, each with a percentile 95\% interval from 200 draws of spatial blocks of the slide. In place: the nuclei where the registration puts them, median over nuclei. Background: the same statistic with every nucleus first moved 25 capture locations (50\,\textmu m) on each axis, along each of the four diagonals, median over the four placements of every nucleus; a moved nucleus can overlap other nuclei, so this is a background placement, not a nucleus-free one. Nuclei: nuclei in place after leaving out those that would cross the edge of the capture area at one of these shifts and those whose mean over these shifts is zero. Deposited map: the nucleus map deposited with the dataset. Our Cellpose: our masks of the ovarian dataset moved by ($-$3, $-$3) locations, the dip under our registration. Xenium: the nuclei imaged by Xenium on the same section as the Visium HD run, in four cores of one slide.}
\label{supptab:7}
\setlength{\tabcolsep}{4pt}
\begin{tabular}{llrrrrr}
\toprule
Case & Nuclei from & Nuclei & Against & In place & Background & In place $-$ background \\
\midrule
Ovarian & deposited map & 216,104 & neighbours & 4.00 [3.94, 4.09] & 0.63 [0.60, 0.65] & 3.37 [3.31, 3.46] \\
 &  &  & next out & 9.62 [9.45, 9.79] & 1.87 [1.83, 1.92] & 7.75 [7.58, 7.91] \\
\addlinespace[2pt]
Ovarian & our Cellpose & 160,309 & neighbours & 3.04 [2.97, 3.12] & 0.46 [0.44, 0.48] & 2.59 [2.52, 2.67] \\
 &  &  & next out & 8.03 [7.84, 8.21] & 1.46 [1.41, 1.50] & 6.58 [6.39, 6.74] \\
\addlinespace[2pt]
Colorectal & deposited map & 273,394 & neighbours & 4.57 [4.41, 4.74] & 0.87 [0.83, 0.90] & 3.70 [3.55, 3.86] \\
 &  &  & next out & 9.45 [9.19, 9.73] & 2.45 [2.39, 2.53] & 7.00 [6.74, 7.24] \\
\addlinespace[2pt]
Colorectal, H1-VM2JXXK/A1 & deposited map & 192,523 & neighbours & 5.58 [5.44, 5.72] & 1.42 [1.37, 1.49] & 4.15 [4.01, 4.33] \\
 &  &  & next out & 9.88 [9.61, 10.17] & 3.52 [3.38, 3.64] & 6.36 [6.05, 6.74] \\
\addlinespace[2pt]
Lung core VUILD49LA & Xenium & 32,271 & neighbours & 5.38 [4.76, 5.88] & 3.11 [3.03, 3.45] & 2.26 [1.60, 2.85] \\
 &  &  & next out & 9.43 [8.57, 10.36] & 8.05 [7.69, 8.57] & 1.39 [0.46, 2.58] \\
\addlinespace[2pt]
Lung core TILD113LA & Xenium & 18,178 & neighbours & 5.08 [4.59, 5.68] & 3.03 [2.54, 3.19] & 2.05 [1.56, 2.94] \\
 &  &  & next out & 9.43 [8.43, 10.62] & 7.69 [7.25, 8.20] & 1.74 [0.53, 2.99] \\
\addlinespace[2pt]
Lung core TILD111LA & Xenium & 22,708 & neighbours & 3.70 [3.03, 4.26] & 3.03 [3.03, 3.45] & 0.67 [$-$0.21, 1.20] \\
 &  &  & next out & 6.36 [5.34, 7.34] & 7.69 [7.27, 8.16] & $-$1.34 [$-$2.42, $-$0.33] \\
\addlinespace[2pt]
Lung core VUHD049 & Xenium & 19,113 & neighbours & 2.61 [2.04, 3.03] & 4.00 [3.45, 4.00] & $-$1.39 [$-$1.96, $-$0.58] \\
 &  &  & next out & 4.99 [4.00, 5.88] & 8.99 [8.57, 9.44] & $-$4.01 [$-$5.09, $-$2.69] \\
\bottomrule
\end{tabular}
\end{table}
}

\supptable{8}{%
\begin{table}[H]
\centering\scriptsize
\caption{\textbf{Supplemental Table S8. Whole-slide and regional offsets under each registration: the dip and the correlation ratio.} Our nucleus masks were placed on the capture grid by each registration and moved over it one capture location at a time. Shifts are (rows, columns) of capture locations of 2\,\textmu m, rows running down the slide. Dip: the position of the deepest dip within 12 locations; Depth: one minus the mean count at the dip divided by the median over its 8 neighbouring shifts. Correlation ratio: the capture locations are grouped by their distance from the edge of the nearest nucleus, in bands one location wide (3 inside the nuclei, 6 outside them, and the rest of the tissue); each count is divided by the mean count of the capture locations in tissue in the 5 $\times$ 5 square around it; the shift is the highest local maximum of the correlation ratio $\eta^2$ of the divided counts on the groups within 12 locations, higher than at its 8 neighbouring shifts (Supplemental Methods). Offset: the distance of the shift from zero. Cellpose: nuclei segmented with Cellpose-SAM, as \STARgrid{} does; StarDist: nuclei segmented with StarDist, the control on the ovarian dataset and the two colorectal slides. Regions: the capture area divided 5 $\times$ 5; a region is measured when it holds at least 30,000 nucleus locations in tissue. Usable: not on the edge of the search, and in at least 95\% of 200 draws of spatial blocks of the slide the shift found lies within one location of the full-data one. At zero: usable regions whose shift is zero. Largest: the largest regional offset. Both usable: of the regions in which the dip and the correlation ratio are both usable, those where they give the same shift (Equal) and shifts one location apart (One off); in none do they differ by more. Deposited: the registration in the output deposited with the dataset; our run: our Space Ranger run on a subsample of 100,000 reads (Supplemental Methods). For each set of nuclei on a slide, the masks and the count image are the same in every row; only the registration changes. Every whole-slide dip and correlation ratio is usable, and on every whole slide the two give the same shift.}
\label{supptab:8}
\setlength{\tabcolsep}{3pt}
\begin{tabular}{llrrrrrrrrrr}
\toprule
 & & \multicolumn{2}{c}{Whole slide, dip} & \multicolumn{2}{c}{Whole slide, ratio} & \multicolumn{4}{c}{Regions usable} & \multicolumn{2}{c}{Both usable} \\
\cmidrule(lr){3-4}\cmidrule(lr){5-6}\cmidrule(lr){7-10}\cmidrule(l){11-12}
Nuclei & Registration & Shift & Depth (\%) & Shift & Offset (\textmu m) & Dip & Ratio & At zero & Largest (\textmu m) & Equal & One off \\
\midrule
\multicolumn{12}{l}{\textit{Ovarian dataset, H1-YQJBZ7X/D1}} \\
Cellpose & Ours, from the images & ($-$3, $-$3) & 2.3 & ($-$3, $-$3) & 8.5 & 21 & 21 & 0 & 10.0 & 20 & 1 \\
StarDist & Ours, from the images & ($-$3, $-$3) & 3.0 & ($-$3, $-$3) & 8.5 & 21 & 21 & 0 & 10.0 & 20 & 1 \\
Cellpose & Space Ranger 4.0.1, deposited & (0, 0) & 3.0 & (0, 0) & 0.0 & 21 & 21 & 20 & 2.0 & 20 & 1 \\
StarDist & Space Ranger 4.0.1, deposited & (0, 0) & 4.0 & (0, 0) & 0.0 & 20 & 20 & 20 & 0.0 & 20 & 0 \\
\addlinespace[3pt]
\multicolumn{12}{l}{\textit{Colorectal slide, H1-GMHFWPH/D1}} \\
Cellpose & Ours, from the images & (0, 0) & 3.1 & (0, 0) & 0.0 & 21 & 21 & 13 & 2.0 & 19 & 2 \\
StarDist & Ours, from the images & (0, 0) & 3.5 & (0, 0) & 0.0 & 21 & 20 & 13 & 2.0 & 18 & 2 \\
Cellpose & Space Ranger 4.1.0, deposited & (0, 0) & 3.3 & (0, 0) & 0.0 & 21 & 21 & 21 & 0.0 & 20 & 1 \\
StarDist & Space Ranger 4.1.0, deposited & (0, 0) & 3.8 & (0, 0) & 0.0 & 21 & 20 & 20 & 0.0 & 20 & 0 \\
\addlinespace[3pt]
\multicolumn{12}{l}{\textit{SPATCH slide, H1-WP2FTZD/D1}} \\
Cellpose & Ours, from the images & ($-$5, $-$3) & 2.0 & ($-$5, $-$3) & 11.7 & 22 & 22 & 0 & 15.2 & 20 & 2 \\
Cellpose & Ours, with our fitted correction & (0, 0) & 4.1 & (0, 0) & 0.0 & 22 & 22 & 22 & 0.0 & 21 & 1 \\
Cellpose & Space Ranger 3.0.0, our run & ($-$5, $-$2) & 2.4 & ($-$5, $-$2) & 10.8 & 22 & 22 & 0 & 12.6 & 18 & 4 \\
Cellpose & Space Ranger 4.1.0, our run & (0, 0) & 2.4 & (0, 0) & 0.0 & 22 & 22 & 7 & 4.5 & 18 & 4 \\
\addlinespace[3pt]
\multicolumn{12}{l}{\textit{Colorectal slide H1-VM2JXXK/A1}} \\
Cellpose & Space Ranger 3.0.0, deposited in GEO & ($-$1, $-$2) & 2.7 & ($-$1, $-$2) & 4.5 & 13 & 16 & 0 & 6.3 & 11 & 2 \\
StarDist & Space Ranger 3.0.0, deposited in GEO & ($-$1, $-$2) & 3.1 & ($-$1, $-$2) & 4.5 & 18 & 20 & 0 & 7.2 & 16 & 1 \\
Cellpose & Space Ranger 4.1.0, our run & (0, 0) & 3.0 & (0, 0) & 0.0 & 13 & 16 & 12 & 2.0 & 12 & 1 \\
StarDist & Space Ranger 4.1.0, our run & (0, 0) & 3.8 & (0, 0) & 0.0 & 17 & 20 & 14 & 2.0 & 14 & 2 \\
\bottomrule
\end{tabular}
\end{table}
}

\supptable{11}{%
\begin{table}[H]
\centering\scriptsize
\caption{\textbf{Supplemental Table S11. Lung tissue measured by Xenium and then Visium HD on the same section.} Tissue: the four cores of one slide (H1-K3WKZHN/D1) and the whole sections of two further slides (identifiers given). Each tissue has two rows for the nuclei segmented from the H\&E image: Cellpose-SAM, as \STARgrid{} does, and StarDist, the control; the Xenium nuclei are the same in both rows and their results are given in the first. Registration: Space Ranger 3.0.0 in all of them; by hand for the four cores and, in Loupe Browser 8.0, for H1-JGHRVB7/D1. Mask displacement: the shift that brings the nuclei segmented from the H\&E image onto the nuclei imaged by Xenium, measured from the two masks alone, without counts, for the segmentation of the row; this is the truth for the H\&E nuclei of the row, and zero is the truth for the Xenium nuclei. Dip: the deepest dip within 12 locations in the counts deposited with the dataset. Correlation ratio and Usable: as in Supplemental Table~S8. Within one: the percentage of draws of spatial blocks in which the shift found lies within one location of the full-data one. Shifts not usable are in italics. Shifts are (rows, columns) in capture locations of 2\,\textmu m, rows running down the slide. Lower part, regions (the 3 $\times$ 3 division of each core, the 5 $\times$ 5 division of a section; regions holding at least 30,000 nucleus locations in tissue): the truth of a region is the mask displacement of that region for the H\&E nuclei and zero for the Xenium nuclei. Usable dips and usable ratios: regions where the dip or the correlation ratio is usable, by the distance of its shift from the truth. No usable dip lies two or more locations from its truth.}
\label{supptab:11}
\setlength{\tabcolsep}{2.5pt}
\begin{tabular}{lrrrrrrrrrr}
\toprule
 & \multicolumn{2}{c}{Mask displacement} & \multicolumn{2}{c}{H\&E nuclei, dip} & \multicolumn{2}{c}{H\&E nuclei, ratio} & \multicolumn{2}{c}{Xenium nuclei, dip} & \multicolumn{2}{c}{Xenium nuclei, ratio} \\
\cmidrule(lr){2-3}\cmidrule(lr){4-5}\cmidrule(lr){6-7}\cmidrule(lr){8-9}\cmidrule(l){10-11}
Tissue & Shift & \textmu m & Shift & Within one (\%) & Shift & Within one (\%) & Shift & Within one (\%) & Shift & Within one (\%) \\
\midrule
VUILD49LA, Cellpose-SAM & ($-$2, $-$3) & 7.2 & ($-$2, $-$4) & 100.0 & ($-$2, $-$4) & 100.0 & (0, 0) & 100.0 & (0, 0) & 100.0 \\
VUILD49LA, StarDist (control) & ($-$2, $-$3) & 7.2 & ($-$2, $-$3) & 100.0 & ($-$2, $-$3) & 100.0 & \multicolumn{4}{c}{Xenium nuclei: as in the row above} \\
TILD113LA, Cellpose-SAM & ($-$4, $-$2) & 8.9 & ($-$4, $-$2) & 100.0 & ($-$4, $-$2) & 100.0 & (0, 0) & 100.0 & (0, 0) & 100.0 \\
TILD113LA, StarDist (control) & ($-$4, $-$2) & 8.9 & ($-$4, $-$2) & 100.0 & ($-$4, $-$2) & 100.0 & \multicolumn{4}{c}{Xenium nuclei: as in the row above} \\
TILD111LA, Cellpose-SAM & ($-$3, $-$3) & 8.5 & \textit{($-$3, $-$3)} & 59.0 & ($-$3, $-$3) & 100.0 & \textit{(0, 0)} & 88.5 & (0, 0) & 100.0 \\
TILD111LA, StarDist (control) & ($-$3, $-$3) & 8.5 & ($-$3, $-$3) & 100.0 & ($-$3, $-$3) & 100.0 & \multicolumn{4}{c}{Xenium nuclei: as in the row above} \\
VUHD049, Cellpose-SAM & ($-$5, $-$1) & 10.2 & none & --- & \textit{($-$4, $-$2)} & 50.5 & none & --- & \textit{($-$1, $-$1)} & 68.5 \\
VUHD049, StarDist (control) & ($-$5, $-$1) & 10.2 & ($-$5, $-$1) & 97.0 & ($-$5, $-$1) & 100.0 & \multicolumn{4}{c}{Xenium nuclei: as in the row above} \\
H1-9QZXXNP/D1, Cellpose-SAM & ($-$4, $-$2) & 8.9 & ($-$5, $-$2) & 100.0 & ($-$5, $-$2) & 100.0 & (0, 0) & 100.0 & (0, 0) & 100.0 \\
H1-9QZXXNP/D1, StarDist (control) & ($-$4, $-$2) & 8.9 & \textit{($-$4, $-$2)} & 78.5 & ($-$4, $-$2) & 100.0 & \multicolumn{4}{c}{Xenium nuclei: as in the row above} \\
H1-JGHRVB7/D1, Cellpose-SAM & ($-$4, $-$2) & 8.9 & ($-$4, $-$2) & 100.0 & ($-$4, $-$2) & 100.0 & (0, 0) & 100.0 & (0, 0) & 100.0 \\
H1-JGHRVB7/D1, StarDist (control) & ($-$4, $-$2) & 8.9 & \textit{($-$4, $-$2)} & 65.0 & ($-$4, $-$2) & 100.0 & \multicolumn{4}{c}{Xenium nuclei: as in the row above} \\
\bottomrule
\end{tabular}
\medskip
\begin{tabular}{lrrrrrrrrrr}
\toprule
 & \multicolumn{6}{c}{H\&E nuclei} & \multicolumn{4}{c}{Xenium nuclei} \\
\cmidrule(lr){2-7}\cmidrule(l){8-11}
 & & \multicolumn{2}{c}{Usable dips} & \multicolumn{3}{c}{Usable ratios} & & \multicolumn{3}{c}{Usable ratios} \\
\cmidrule(lr){3-4}\cmidrule(lr){5-7}\cmidrule(l){9-11}
Tissue & Regions & At truth & One off & At truth & One off & Two off & Regions & At truth & One off & Two off \\
\midrule
VUILD49LA, Cellpose-SAM & 0 & 0 & 0 & 0 & 0 & 0 & 5 & 5 & 0 & 0 \\
VUILD49LA, StarDist (control) & 9 & 1 & 2 & 3 & 3 & 0 & \multicolumn{4}{c}{as above} \\
TILD113LA, Cellpose-SAM & 0 & 0 & 0 & 0 & 0 & 0 & 0 & 0 & 0 & 0 \\
TILD113LA, StarDist (control) & 9 & 5 & 3 & 5 & 3 & 0 & \multicolumn{4}{c}{as above} \\
TILD111LA, Cellpose-SAM & 0 & 0 & 0 & 0 & 0 & 0 & 0 & 0 & 0 & 0 \\
TILD111LA, StarDist (control) & 8 & 0 & 3 & 3 & 3 & 0 & \multicolumn{4}{c}{as above} \\
VUHD049, Cellpose-SAM & 0 & 0 & 0 & 0 & 0 & 0 & 0 & 0 & 0 & 0 \\
VUHD049, StarDist (control) & 2 & 0 & 0 & 1 & 0 & 0 & \multicolumn{4}{c}{as above} \\
H1-9QZXXNP/D1, Cellpose-SAM & 17 & 13 & 0 & 14 & 3 & 0 & 17 & 17 & 0 & 0 \\
H1-9QZXXNP/D1, StarDist (control) & 16 & 2 & 0 & 11 & 3 & 0 & \multicolumn{4}{c}{as above} \\
H1-JGHRVB7/D1, Cellpose-SAM & 17 & 7 & 3 & 11 & 4 & 0 & 17 & 17 & 0 & 0 \\
H1-JGHRVB7/D1, StarDist (control) & 17 & 2 & 0 & 8 & 6 & 1 & \multicolumn{4}{c}{as above} \\
\bottomrule
\end{tabular}
\end{table}
}

\supptable{6}{%
\begin{table}[H]
\centering\scriptsize
\caption{\textbf{Supplemental Table S6. Where the molecules STARgrid counts beyond Space Ranger come from, on two slides.} $M$: the molecules that Space Ranger counted (Space Ranger 4.0.1 on slide H1-VM2JXXK/A1, on the genes of its probe set; 4.1.0 with introns included on the ovarian dataset). $T$: STARgrid's total. $B$ and $H$ count the reads that Space Ranger mapped to a single gene and gave a UMI, with Space Ranger's gene and the raw UMI, at Space Ranger's locations ($B$) or at the locations STARgrid chooses ($H$), under the UMI correction of the matrices. Columns of the lower rows: $D=T-M$; placement $H-B$; gene assignment $T-H$ (STARgrid's genes, selection of reads and handling of UMIs with reads for two genes); deduplication $B-M$. The three components add up to $D$.}
\label{supptab:6}
\setlength{\tabcolsep}{5pt}
\begin{tabular}{lrrrr}
\toprule
 & $M$ / $D$ & $T$ / $H-B$ & $B$ / $T-H$ & $H$ / $B-M$ \\
\midrule
Colorectal (Flex), H1-VM2JXXK/A1 & 272,612,275 & 279,777,220 & 273,110,393 & 280,257,120 \\
\quad molecules & +7,164,945 & +7,146,727 & $-$479,900 & +498,118 \\
\quad \% of $D$ & 100 & +99.7 & $-$6.7 & +7.0 \\
\addlinespace[3pt]
Ovarian dataset (RNA-seq), H1-YQJBZ7X/D1 & 187,552,358 & 201,888,275 & 189,504,493 & 198,231,870 \\
\quad molecules & +14,335,917 & +8,727,377 & +3,656,405 & +1,952,135 \\
\quad \% of $D$ & 100 & +60.9 & +25.5 & +13.6 \\
\bottomrule
\end{tabular}
\end{table}
}

\supptable{10}{%
\begin{table}[H]
\centering\scriptsize
\caption{\textbf{Supplemental Table S10. The dip and the correlation ratio on four public Visium HD datasets processed with Space Ranger 4.0.1.} Nuclei: the nucleus map deposited with each dataset, placed by the deposited registration; counts: the deposited 2\,\textmu m counts in tissue. Dip: the deepest dip within 12 locations of the whole slide, in (rows, columns) of capture locations of 2\,\textmu m; depth is one minus the mean count at the dip divided by the median over its 8 neighbouring shifts, in \%. Correlation ratio and Usable: as in Supplemental Table~S8. Regions: of the 5 $\times$ 5 regions holding at least 30,000 nucleus locations in tissue, those where the dip, or the correlation ratio, is usable; At zero and One off: usable regional ratios at zero shift, Space Ranger's registration, and one location from it; none is farther. No dip: of the regions with no dip at all, those where the correlation ratio is usable (---, every region has a dip). Inside, Outside: mean count per capture location in tissue inside nuclei (all locations of nuclei) and 0 to 2\,\textmu m outside the nucleus edge. Every whole-slide correlation ratio is usable.}
\label{supptab:10}
\setlength{\tabcolsep}{3pt}
\begin{tabular}{lrrrrrrrrrrr}
\toprule
 & \multicolumn{3}{c}{Whole slide, dip} & Ratio & \multicolumn{5}{c}{Regions usable} & \multicolumn{2}{c}{Mean count} \\
\cmidrule(lr){2-4}\cmidrule(lr){5-5}\cmidrule(lr){6-10}\cmidrule(l){11-12}
Tissue & Shift & Depth (\%) & Usable & Shift & Dip & Ratio & At zero & One off & No dip & Inside & Outside \\
\midrule
human lymph node FFPE & (0, 0) & 7.1 & yes & (0, 0) & 10 of 10 & 9 of 10 & 7 & 2 & --- & 9.1 & 10.5 \\
human prostate cancer FFPE & (0, 0) & 2.7 & yes & ($-$1, 1) & 14 of 23 & 19 of 23 & 6 & 13 & 5 of 6 & 6.8 & 6.2 \\
human kidney FFPE & no dip & --- & no & (0, 0) & 2 of 22 & 22 of 22 & 15 & 7 & 1 of 1 & 9.8 & 9.0 \\
mouse brain FFPE & no dip & --- & no & (0, $-$1) & 0 of 10 & 10 of 10 & 2 & 8 & 3 of 3 & 30.7 & 24.1 \\
\bottomrule
\end{tabular}
\end{table}
}

\supptable{9}{%
\begin{table}[H]
\centering\scriptsize
\caption{\textbf{Supplemental Table S9. The window of the local mean.} The correlation ratio (Supplemental Methods) with each count divided by the mean count of the capture locations in tissue in a square window around it, of the side given (in capture locations; half-width in \textmu m), or not divided (none); in bold, the window \STARgrid{} uses. Shifts: usable shifts, by the rule of Supplemental Table S8. Lung: the 152 cores, whole sections and regions of Supplemental Table S11, with nuclei from the H\&E image or imaged by Xenium; At truth: shifts at the truth; 2+ off: two or more locations from it. Tumour slides: the 345 whole slides and regions of the placements of Supplemental Table S8; With dip: those where the dip is also usable; 2+ from dip: of these, shifts two or more locations from the dip's. Public slides: the 69 whole slides and regions of Supplemental Table S10; 2+ from SR: shifts two or more locations from Space Ranger 4's registration.}
\label{supptab:9}
\setlength{\tabcolsep}{4pt}
\begin{tabular}{lrrrrrrrrr}
\toprule
 & & \multicolumn{3}{c}{Lung (152)} & \multicolumn{3}{c}{Tumour slides (345)} & \multicolumn{2}{c}{Public slides (69)} \\
\cmidrule(lr){3-5}\cmidrule(lr){6-8}\cmidrule(l){9-10}
Window & Half-width (\textmu m) & Shifts & At truth & 2+ off & Shifts & With dip & 2+ from dip & Shifts & 2+ from SR \\
\midrule
none & --- & 106 & 75 & 4 & 73 & 62 & 10 & 40 & 6 \\
3 $\times$ 3 & 2 & 101 & 81 & 0 & 323 & 310 & 0 & 56 & 3 \\
\textbf{5 $\times$ 5} & \textbf{4} & \textbf{137} & \textbf{109} & \textbf{1} & \textbf{341} & \textbf{328} & \textbf{0} & \textbf{64} & \textbf{0} \\
7 $\times$ 7 & 6 & 136 & 109 & 1 & 343 & 330 & 0 & 66 & 1 \\
9 $\times$ 9 & 8 & 136 & 107 & 1 & 341 & 330 & 0 & 68 & 1 \\
11 $\times$ 11 & 10 & 135 & 108 & 1 & 339 & 328 & 0 & 69 & 1 \\
15 $\times$ 15 & 14 & 131 & 101 & 2 & 320 & 310 & 0 & 69 & 1 \\
21 $\times$ 21 & 20 & 131 & 102 & 1 & 300 & 290 & 3 & 68 & 1 \\
31 $\times$ 31 & 30 & 131 & 100 & 0 & 253 & 242 & 12 & 67 & 1 \\
41 $\times$ 41 & 40 & 131 & 98 & 1 & 199 & 186 & 17 & 64 & 1 \\
61 $\times$ 61 & 60 & 130 & 95 & 1 & 171 & 158 & 22 & 61 & 2 \\
\bottomrule
\end{tabular}
\end{table}
}
